% Vietnamese translation for OI Public Library
% Source-File: IOI中国国家候选队论文/2004/朱泽园.ppt
\documentclass[11pt]{article}
\usepackage{oipl-vn}

\newcommand{\Oh}{\mathcal{O}}

\title{Khớp nhiều xâu và các gợi ý thuật toán\\{\large Ghi chú theo slide}}
\author{Zhu Zeyuan}
\date{IOI 2004}

\begin{document}
\maketitle

\origtitle{Multiple string matching algorithms and their implications}
\sourcepath{IOI 2004 / Zhu Zeyuan.ppt}

\section{Từ KMP đến cấu trúc cây}

Các slide mở bằng các mẫu như \texttt{abcd}, \texttt{bcde} và nhắc
Knuth--Morris--Pratt cùng con trỏ tiền tố. Với một mẫu đơn, KMP cho ta cách
dịch mẫu mà không quét lại các ký tự đã biết. Với nhiều mẫu, cùng ý tưởng
được nâng lên thành trie và các liên kết thất bại: trạng thái hiện tại chính
là tiền tố dài nhất còn có thể khớp.

\section{Hai cây và thao tác quét}

Phần trích được từ slide nhắc \texttt{TreeA}, \texttt{TreeB}, các thủ tục
\texttt{ScanA(Left,Right,P)} và \texttt{ScanB(Left,Right)}. Đây là phần nói
về cách phối hợp cây tiền tố và cây hậu tố để xử lý nhiều xâu, trong đó các
đoạn \(T[\mathit{Left}..\mathit{Right}]\) được quét và dịch bằng mảng
\(\mathit{Shift}\).

Ví dụ \texttt{abcabcde} minh họa việc chuyển từ một đoạn đã khớp sang đoạn
kế tiếp mà không làm lại toàn bộ công việc. Khi các liên kết và phép quét
được tổ chức đúng, mỗi ký tự chỉ gây ra số thao tác hữu hạn, dẫn tới thời
gian tuyến tính trong độ dài văn bản đối với phần lõi.

\section{Ý nghĩa}

Điểm chính của bài trình bày không chỉ là một cấu trúc dữ liệu cụ thể. Nó cho
thấy nhiều thuật toán xâu mạnh đều cùng dựa trên việc ghi nhớ biên khớp:
tiền tố, hậu tố, trạng thái thất bại và liên kết hậu tố. Khi phần thông tin
đó được lưu trong cây, ta xử lý được cả họ mẫu thay vì một mẫu đơn.

\end{document}
